% !TeX root = ../教材题库_全部条目_转写工作稿.tex
% Source: XJ-B1-C03-TH-P083-Y0867
% 第3章 / 3.2 函数的基本性质 / 教材 p.83 / PDF 第90页
% OCR未校对：公式与文字均可能有误；未完成数学TeX转写；禁止直接发给学生
% 原题及相邻原解截图（核对用）：../题目原图/XJ-B1-C03-TH-P083-Y0867-01.webp
% 本文件只在转写完成后作为题干使用；OCR 条目尚不是可靠题干。
\begin{bankitem}{p.83 思考（含紧邻说明）}{第3章 / 3.2 函数的基本性质 / 教材 p.83 / PDF 第90页}{XJ-B1-C03-TH-P083-Y0867}{OCR未校对：公式与文字均可能有误；未完成数学TeX转写；禁止直接发给学生}
\begin{ocrdraft}
[教材p.83 / PDF第90页]
\par
？思考如何用数学符号语言来描述函数的奇偶性？
\par
偶函数F(x）的图象是以y轴为对称轴的轴对称图形，也就是说，若点，A在
\par
F(x）的图象上，则点A关于3，轴的对称点B也在该图象上，而3，轴是线段AB的垂
\par
直平分线．
\par
于是，若点A的坐标为0，F(x))，则点B的坐标为（一产F(—x))；直线AB
\par
平行于轴，所以点A与点B均在x轴同侧，并且到工轴的距离相等，即F(—x)=
\par
反过来，如果对于F0）的定义域中的任意F（一到都有定义并且满足
\par
F(—x)=F(x)，则F(x）的图象上任一点，F0））关于轴的对称点（一
\par
F(x））=（亠r，F（一）也在该图象上，这表明F(x）的图象是以3，轴为对称轴的轴
\par
对称图形．
\par
因此偶函数就是满足条件F（亠）一F(x）的函数．
\par
同理可证：奇函数是满足条件F（亠）一一F(x）的函数．
\par
上面的讨论概括如下：
\par
奇偶性是甬数在定义
\par
域内的整体性质，判断函
\par
（1）如果对一切使F(x）有定义的F（一到也有定义，
\par
数的奇偶性应先明确甬数
\par
并且F（一一F(x），则称F(x）为偶函数；
\par
的定义域．
\par
（2）如果对一切使F(x）有定义的F（一到也有定义，
\par
并且F（一到一一F(x），则称F(x）为奇函数．
\end{ocrdraft}
\end{bankitem}
